% 第 36 章　偏振与光在物质中的传播
\chapter{偏振与光在物质中的传播}

\step{反常识开场}

请你想象一个净水器。水先流过一层滤芯，再流过第二层。如果你在两滤芯之间再加一层更密的滤芯，流出的水只会更少，绝不可能更多。多层过滤只会挡得更严——这是每个人对“过滤”的直觉。

现在把净水器换成光。找一副偏振太阳镜（市面上大多数太阳镜就是），再找一个旧液晶显示器或一部液晶屏手机。液晶屏发出的光，天生就是“整理过的”光——这一点本章会解释。把偏振镜片放在屏幕前，慢慢旋转它：在某个角度，屏幕看起来正常；转过九十度，屏幕几乎全黑。这说明屏幕上来的光被镜片“过滤”掉了。

好，现在做那个不可能的动作。保持第一片镜片不动（屏幕全黑的角度），在它前面再放一片偏振镜，旋转这第二片。你会发现：加到两片之后，无论怎么转，屏幕始终是黑的。这很合理，两层过滤嘛。

接下来是见证奇迹的时刻：把第二片镜片\textbf{斜插在两者之间}，让三片的朝向依次为竖直、斜四十五度、水平。原本漆黑的屏幕上，竟然有光透了出来！挡住光的滤芯，多加一片，光反而通过了。

这不是魔术道具，任何两片偏振镜都能重复这个实验。它直接扇了“过滤”直觉一记耳光：偏振片对光做的事，显然不是简单的“拦住一部分”。本章的任务，就是搞清楚光身上到底还藏着哪一个我们从未注意的“旋钮”，以及这个旋钮如何解释蓝天、彩虹和液晶屏幕。

顺便说一句公平的话：你会觉得这个现象反常识，恰恰说明你的常识很健康。过滤直觉来自水和空气里的一切宏观过滤器，它们对“被滤对象”的认识确实只有“挡多挡少”一个维度。问题在于光不是一种待滤的杂质流——实验马上要证明，它随身携带的那个隐藏维度，宏观过滤器从来没有机会暴露给你看。直觉这次不是错了，而是从没上过班。

\step{先猜一猜}

在往下读之前，请先写下你自己的猜测。下面几种说法，哪一个最像真相？

第一种：偏振片是有方向性的颜色滤镜，变黑是染色叠加深了，第三片恰好把颜色“调”了回来。

第二种：光是沿直线飞的小颗粒流，偏振片是栅栏，颗粒排队从缝里钻过去；三片栅栏的缝碰巧排成一条通道。

第三种：光除了“往哪飞”“什么颜色”“有多亮”之外，还有一个额外的属性，偏振片会读写这个属性；第三片不是“放行”，而是先把这个属性\textbf{改写}了，后面的片才能透过。

写下你的选择。本章结尾我们回来对答案。也请记住净水器直觉给你的预测：“多挡一层，透过更少”。如果这个预测被实验击穿，被击穿的就不只是对偏振片的看法，而是“光可以被完全描述为一束沿轨道飞行的东西”这个更深的图像。

再给一点下注前的提示，但不给答案。第一种说法有一个硬伤：颜色滤镜叠得再多，也不会让全黑的两层因为再加一层而变亮——染色只会减光，不会生光。第二种说法听上去很有画面感，请你特别记住它，因为它正是几百年前一批最聪明的头脑押过的注，而它的命运将在本章中段被正式宣判。第三种说法目前还只是空话——“额外的属性”是什么属性？“改写”又是什么意思？空话要变成物理，得先过实验这一关。我们这就去做实验。

\step{动手实验}

\begin{experiment}[三片“滤芯”与屏幕里的暗格]
材料：一部手机（液晶屏即可，OLED 屏效果差一些）、一副偏振太阳镜。如果找不到第二片偏振镜，可以把 3D 电影眼镜、另一副偏光墨镜，甚至旧计算器上拆下来的偏振片拿来用。

第一步：打开手机的纯白色画面（备忘录空白页就行）。把偏振镜片贴在屏幕前，慢慢旋转一整圈。你会看到亮度经历两次最亮、两次最暗。在最暗的角度停住——此时镜片与屏幕的“方向”正好互相垂直。

第二步：在镜片与屏幕之间，斜着插入第二片偏振镜（与两片都成大约四十五度角）。原本全黑的画面会重新亮起来，大约恢复到不隔镜片时的八分之一亮度。抽出中间那片，立刻变黑；插回去，又亮。这个“插片开光”的动作请多做几次，让惊讶沉淀成问题。

第三步：带着镜片去户外。抬头看与太阳成九十度角方向的天空，旋转镜片，你会看到那片天空明显变暗又变亮——天空的光也是被“整理过”的。再低头看水面或柏油路面上的刺眼反光，旋转镜片到某个角度，反光几乎消失，水下和路面下的东西显露出来。反光消得最干净时，太阳、水面和你眼睛连成的角度大约在五十多度，这个角度本章后面会有名字。
\end{experiment}

实验做完了，三个观察请记在纸上：镜片旋转一周，亮度变化两次而不是一次；中间插片让全黑变亮；水面反光可以在某个特定角度几乎被消灭。这三件事，每一件都是线索。

先做一点热身推理。为什么旋转一周是两次最亮、两次最暗，而不是一次？如果偏振镜是在认“方向”，那么转到一百八十度时，镜片的状态其实和出发时一模一样（把一片玻璃旋转半圈，它没有变得“上下颠倒”——它没有上下），所以亮度必然在转半圈后就回到原样。这个小小的对称性论证提示我们：光身上那个隐藏的旋钮，是一种“轴”的性质，而不是一种“箭头”的性质——轴转半圈复原，箭头要转一整圈。先把这个区别存在心里，它后面会多次出现。

\step{旧模型遇到麻烦}

我们清点一下手里的旧模型，看它们各自在哪里卡壳。

第 33 章把光当成射线：一条光线由位置和方向完全决定。这个模型能解释影子、反射、折射，但它给光线安排的“档案”里根本没有“绕轴转了九十度”这一栏。一条沿直线飞行的线，绕着自己的轴旋转，还是那条线。所以射线模型预言：旋转偏振片什么都不会发生。实验说它错了。

第 35 章把光当成波。进步很大，但还不够：波有两种。像声波那样的纵波，振动方向与传播方向一致；像绳波那样的横波，振动方向与传播方向垂直。如果光是纵波，它同样没有“绕轴方向”可言——纵波绕着自己的传播轴旋转，也还是原来的波，正如一根前后伸缩的弹簧怎么转都一样。声波穿过空气时，无论扬声器横放竖放，你听到的都一样。所以“光是纵波”也解释不了旋转镜片的变化。

那就只剩一条路：光是横波，而且不同的横波，振动方向可以不同。一个沿水平方向传播的横波，可以上下振，可以左右振，也可以斜着振。绕传播轴旋转，真的会换到另一个状态。可是新的麻烦立刻出现：自然光（阳光、灯光）穿过一片偏振片，无论镜片转到什么角度，出来的亮度都一样，而且都恰好减半。如果自然光也有确定的振动方向，为什么不挑角度？更麻烦的是第三步观察：透过一片的光再经过垂直的第二片，全灭；中间斜插一片，又活了。“过滤”类模型对此完全无能为力——过滤器只能删，不能改。

\begin{mainline}
旧模型遇到的麻烦可以压成一句话：\textbf{光除了飞行方向、颜色和亮度，还携带一个“绕轴的方向”，自然光是这个方向的一锅乱炖，而偏振片不是筛子——它是一台能把光的这个方向投影、改写的机器。}我们需要一个专门的名字来描述这个方向。
\end{mainline}

\step{发明新概念}

现在我们自己动手造概念。先做一个机械模型找感觉：拿一根长绳，一端系在门把手上，你拿着另一端上下甩，绳上就跑起一列上下振动的横波。再左右甩，就得到一列左右振动的横波。上下和左右，是同一根绳上两种独立的“振动款式”。把绳子穿过一道竖直的栅栏缝：上下振动的波能通过缝（绳子在缝里照样上下动），左右振动的波被缝卡死。对绳波来说，栅栏真的是筛子。

光比绳波深一层。电磁学告诉我们：光是电磁波，空间上和时间上变化的电场与磁场共同满足麦克斯韦方程、以波的形式向前传。电场振荡的那个方向，就是光的\textbf{偏振方向}。它像绳波的“振动款式”，因为同样是横波；它不像绳波的地方更关键——绳波振的是看得见摸得着的绳子，光振的是电场本身，不需要任何介质。

在分门别类之前，先建立一个后面要反复使用的基本功：任何方向的线振动，都可以拆成上下和左右两份。这不是光学的新定律，只是向量分解——一根斜四十五度的绳子，你让它斜着振，等价于它同时上下振一份、左右振一份，两份一样大、节拍一致。反过来，上下一份加左右一份，合起来就是斜向振动。这套“拆与合”的语言，就是偏振片“改写”偏振方向的全部秘密：偏振片先把你给它的振动拆成沿透光轴和垂直透光轴两份，拦下垂直那份，放行走沿轴那份——出来的光偏振方向就换成了透光轴的方向。下面我们把光的偏振分门别类：

\textbf{线偏振}：电场始终沿一条固定直线来回振荡。绳波上下振、左右振都是线偏振的亲戚。透过理想偏振片的光就是线偏振光，偏振方向由偏振片的“透光轴”决定。

\textbf{自然光（非偏振光）}：阳光和灯光来自亿万个原子各自独立的辐射，每一小段的偏振方向都是随机的，混在一起，任何方向上都不占多数。它绕轴完全对称，所以第一片偏振镜转到哪个角度都一样——无论怎么量，平均值都是一半。这就是“过一片减半”的原因。

\textbf{圆偏振}：想象电场的大小不变，方向却像螺旋桨一样匀速旋转，光向前飞，电场箭头沿途画出一根弹簧的形状。这听上去像一种全新的妖魔鬼怪，其实它有一个朴素的分解：把上下振动和左右振动各拿一份，让两者错开四分之一个节拍（一个到最高点时另一个正好过零点），合起来就是旋转。电影院发的 3D 眼镜用的就是它。

\textbf{椭圆偏振}：圆偏振的更一般情形——旋转的同时电场大小也在变，箭头扫出一个椭圆。线偏振和圆偏振都是它的特例：压扁到一根线是线偏振，圆成正圆是圆偏振。

还有介于中间的\textbf{部分偏振光}：比如天空的光，一大半偏振、一小半乱炖，所以镜片只能让它变暗，不能让它全黑。

圆偏振不只是数学好奇心，它已经坐在你的鼻梁上。现代 3D 电影的两台放映机，一台放左旋圆偏振、一台放右旋圆偏振；眼镜的左镜片只放行左旋、右镜片只放行右旋，于是左右眼各看各的画面，大脑把两幅图拼成立体。用圆偏振而不用线偏振的原因很实际：你歪头时，“竖直”和“水平”会跟着歪，左右眼画面就串了；而“顺时针转”和“逆时针转”不管头怎么歪都不会混淆。上一代用红蓝滤色片的老 3D 眼镜，滤的是颜色；这一代滤的是旋转方向——滤的对象从“光的哪个属性”上升级了。

自然界比工程师更早用上偏振。蜜蜂的复眼能感知天空光的偏振方向，而晴天头顶的偏振图案由太阳位置唯一决定：与太阳成九十度角方向的天空偏振最强，图案像一张以太阳为中心的方向地图。蜜蜂靠这张地图在阴天、在太阳被山挡住时照样导航回巢。你在实验第三步里用镜片看到的“天空变暗变亮”，就是这张蜜蜂地图的边缘。

有了偏振这个概念，三片偏振片的奇迹已经露出解法的一半：中间那片不是“又过滤了一次”，而是把竖直偏振的光\textbf{改写}成斜四十五度偏振的光——斜四十五度的振动里，含有水平方向的分量，于是水平的第三片又有东西可放行了。解法剩下的另一半，要回答“改写多少”，那是数学翻译器的活儿。

概念清单还差最后一块：本章标题的后一半——“光在物质中的传播”。偏振之所以重要，一半原因是它直接参与了光与物质的每一次交手。物质由带电的粒子组成：带负电的电子、带正电的原子核。光波的电场扫过物质时，会推动这些电荷来回振荡；振荡的电荷自己也是小小的辐射源，会向四周重新发出电磁波。光在物质里的全部故事——折射、色散、吸收、散射、双折射——都是“光驱动电荷、电荷再辐射”这一出戏的不同演法：

\textbf{折射率的来源}：在玻璃、水这类透明物质里，入射光与电荷再辐射的光叠加后，波峰的位置被推迟了，整体效果相当于光在物质里走得慢了。第 33 章里那个折射率 \(n\)，微观上就是这场“推迟”的度量。

\textbf{色散}：物质里的电子不是完全自由的，它们被原子核牵着，像拴在弹簧上的小球，有自己偏好的固有频率。入射光的频率离固有频率越近，电荷被带动的幅度越大，推迟得越厉害。紫光的频率比红光更接近玻璃中电子的固有频率，所以紫光被推迟得更多，折射率更大，折得更弯。棱镜把白光分成彩虹，就是因为玻璃对不同颜色“慢待”的程度不同。

\textbf{吸收}：如果光的频率恰好撞上电荷的固有频率，就像推秋千推在点子上，电荷被剧烈驱动，并把能量通过碰撞传给周围的邻居——光的能量真正变成了物质的内能，这就是吸收。玻璃对可见光的固有频率都在紫外区，可见光推不动它们，所以玻璃透明。

\textbf{散射}：空气中的分子也会把光能量抓过来再撒向四面八方。分子比波长小得多时，它们辐射的本事随频率急剧上升——蓝紫光被撒得远比红光凶。你抬头看到的蓝天，就是被空气分子撒向各个方向、其中格外多的蓝光；而日落时太阳本身显得红，是因为蓝光大半在路上被撒掉了，剩下的直奔你眼睛。

\textbf{双折射}：以上默认物质各个方向性质相同。方解石这类晶体内部，不同方向上的“弹簧”松紧不同——光沿某个偏振方向走时遇到的阻力小、沿垂直方向走时阻力大。于是同一块晶体对不同偏振方向有两个折射率，一束光进去被拆成两束。把方解石压在纸上的一个字上，你能看到两个错开的字。偏振在这里从“观察对象”变成了“探针”：它测出了物质内部的方向性。

还有一类物质把“驱动与再辐射”演到了极致：金属。金属里的电子几乎不被原子核拴住，是自由电子，对任何频率的光都能立刻响应，而且响应得太好——它们集体振荡产生的再辐射，在金属内部恰好把入射光抵消干净，同时把一份强烈的反射波送回空气。所以金属不透明却亮闪闪。同一条“电荷接力”定律，电子被拴住时演出透明和折射，电子自由时演出反光和金属光泽。

\step{一句话模型}

\begin{oneline}
光是横波，电场振荡的方向叫偏振；偏振片不删光，它把光的偏振投影到自己的透光轴上；而光进入物质后的一切行为——变慢、分色、被吸收、被撒向天空——都来自电场驱动电荷、电荷再辐射时的大小与快慢之别。
\end{oneline}

这句话请拆开读两遍。第一遍读“投影”：偏振片的动作不是减法，是把电场向量沿透光轴取分量，这正是三片偏振片能“复活”光线的机关。第二遍读“驱动与再辐射”：光与物质的相互作用不是光“钻进”物质找路，而是一场电荷之间的接力，接力的快慢由频率和物质结构共同决定。

\step{数学翻译器}

直觉已经到位，现在把它翻译成可以算的式子。

先看一片偏振片。设入射线偏振光的电场振幅为 \(E_0\)，偏振方向与透光轴夹角为 \(\theta\)。向量的投影是初中几何：沿透光轴的分量大小是 \(E_0\cos\theta\)，垂直分量被吸收。而光的强度正比于电场振幅的平方（这一点与“绳子甩得越猛能量越大”的平方关系同源），所以透射强度
\[
  I = I_0 \cos^2\theta .
\]
这就是\textbf{马吕斯定律}。式子左边回答“出来多亮”，右边每一项都在改变什么：\(I_0\) 是进来的本钱，\(\cos^2\theta\) 是投影两次的结果（先投影振幅，再把振幅平方成强度）。

\begin{center}
\begin{tikzpicture}[scale=1.7,>=Stealth]
  \draw[->] (0,0) -- (2.8,0) node[right] {透光轴};
  \draw[->] (0,0) -- (0,1.8);
  \draw[->,very thick] (0,0) -- (40:2.3) node[above right=-2pt] {$E_0$};
  \draw[->,very thick,blue!60!black] (0,0) -- ({2.3*cos(40)},0) node[below] {$E_0\cos\theta$};
  \draw[dashed] (40:2.3) -- ({2.3*cos(40)},0);
  \draw (0.8,0) arc[start angle=0,end angle=40,radius=0.8];
  \node at (20:1.1) {$\theta$};
\end{tikzpicture}
\end{center}

按规矩，立刻用特殊情形检查它。\(\theta=0\)：\(\cos^2\theta=1\)，全通过——方向已经对齐，没什么可投的，合理。\(\theta=90^\circ\)：\(\cos^2\theta=0\)，全灭——垂直方向的水平分量就是零，合理，这就是实验第一步的全黑。\(\theta=45^\circ\)：\(\cos^2\theta=1/2\)，过一半。再检查自然光：各种方向均等，\(\cos^2\theta\) 对所有方向平均恰好是 \(1/2\)，所以自然光过任何一片都减半，与实验相符。

现在算那个奇迹。自然光强度 \(I_0\) 过第一片：\(I_0/2\)，变成竖直偏振。中间片斜 \(45^\circ\)：\((I_0/2)\times\cos^2 45^\circ = I_0/4\)，变成斜向偏振。第三片水平，与斜向又差 \(45^\circ\)：\((I_0/4)\times\cos^2 45^\circ = I_0/8\)。抽掉中间片，竖直对水平，\(I_0/2\times\cos^2 90^\circ=0\)。八分之一的复活，与动手实验里“大约八分之一的亮度”对上了。净水器直觉正式阵亡：偏振片的级联不是乘法过滤器，而是逐片的旋转接力。

再来一个带真实数字的估算，看看 \(1/\lambda^4\) 在日落时干了什么。地球半径约 \(6400\,\mathrm{km}\)，而大气的有效厚度只有约 \(8\,\mathrm{km}\)。正午的阳光几乎是竖直穿过这 \(8\,\mathrm{km}\)；日落时阳光贴着地面切进来，在同样的空气里要走的路程约是竖直方向的几十倍（粗略的几何就能给出三十倍以上）。散射每走一段路就按波长四次方的偏好“抽税”一次：蓝光的税率是红光的四倍多。走几十倍的路程，等于这税被连续抽了几十次，蓝光几乎被抽干，红光还剩不少。所以正午的太阳白得耀眼，日落的太阳红得温柔——同一个太阳，同一片空气，差别只在光路的长度被几何放大了一个数量级。

\begin{mathbridge}
动手实验里水面反光消失的角度，也可以算出来。光在两种介质界面反射时，平行于入射面偏振的分量和垂直于入射面偏振的分量，反射本领不同。存在一个特殊角度，使平行分量的反射系数恰好为零：此时反射光与折射光的方向互相垂直，再辐射的电荷振荡方向正对着反射方向——而横波不能沿自己的振荡方向传播，于是那个方向的反射干脆消失了。由折射定律和这个垂直条件可得
\[
  \tan\theta_{\mathrm B} = \frac{n_2}{n_1}.
\]
这个角叫\textbf{布儒斯特角}。代入水的折射率 \(n_2=1.33\)、空气 \(n_1=1\)：\(\theta_{\mathrm B}=\arctan 1.33\approx 53^\circ\)，正是实验里反光最弱的“五十多度”。玻璃取 \(n_2=1.5\)，得约 \(56^\circ\)。

\begin{center}
\begin{tikzpicture}[scale=1.1,>=Stealth]
  \draw[thick] (-4,0) -- (4,0) node[right] {水面};
  \fill[blue!8] (-4,-1.6) rectangle (4,0);
  \draw[dashed] (0,0) -- (0,2.6);
  \draw[->,very thick] ({-2.6*tan(53)},2.6) -- (0,0) node[midway,left,xshift=-2pt] {入射光};
  \draw[->,very thick] (0,0) -- ({2.6*tan(53)},2.6) node[midway,right,xshift=2pt] {反射光（线偏振）};
  \draw[->,very thick,blue!50!black] (0,0) -- ({1.9*tan(37)},-1.9) node[midway,right,xshift=2pt] {折射光};
  \draw (0,1.1) arc[start angle=90,end angle=143,radius=1.1];
  \node at (118:1.55) {$\theta_{\mathrm B}$};
  \draw (37:0.7) arc[start angle=37,end angle=-53,radius=0.7];
  \node at (-10:1.05) {$90^\circ$};
  \node[align=center] at (0,-2.3) {布儒斯特角入射时，反射光与折射光恰好互相垂直};
\end{tikzpicture}
\end{center}

检查极端情形：\(n_2=n_1\) 时 \(\theta_{\mathrm B}=45^\circ\)，但此时界面根本不存在，“全偏振反射”无从谈起——公式给的角再漂亮，也得有界面才有意义。
\end{mathbridge}

\begin{mathbridge}
色散与散射也各有一句可以算的。玻璃中电子的固有频率在紫外，可见光频率 \(\omega\) 比它低时，折射率随频率单调缓升，所以紫光的折射率大于红光，彩虹的颜色顺序由此固定。而比波长小得多的粒子，散射强度正比于频率的四次方，也就是反比于波长四次方：
\[
  I_{\text{散}} \propto \frac{1}{\lambda^4}.
\]
取蓝光 \(\lambda\approx 450\,\mathrm{nm}\)、红光 \(\lambda\approx 650\,\mathrm{nm}\)，散射本领之比为 \((650/450)^4\approx 4.3\)——蓝光被撒向天空的效率是红光的四倍多。这就是“蓝天不紫”的数量级起点（还要再乘上太阳光谱和人眼敏感度，紫色才最终没赢过蓝色）。
\end{mathbridge}

\begin{challenge}
线偏振、圆偏振、椭圆偏振可以用一个统一的记号装起来：把电场沿两个垂直方向的分量写成两个数，再配上两者的相位差，组成一个二元对象（琼斯向量）。水平线偏振是 \((1,0)\)，竖直是 \((0,1)\)，斜四十五度是 \((1,1)/\sqrt{2}\)，圆偏振是 \((1,\pm\mathrm{i})/\sqrt{2}\)——虚数单位 \(\mathrm{i}\) 在这里只负责记录“错开四分之一节拍”这个相位差。每一片光学元件（偏振片、双折射晶片）都变成一个 \(2\times2\) 矩阵，整串光路就是矩阵连乘。三片偏振片的 \(I_0/8\)，在这个语言里是三行矩阵乘法的结果。这套记法会在第 46 章以几乎原样的形式回来：量子力学里，偏振是最简单的两态系统，\((1,0)\) 和 \((0,1)\) 会改名叫基态矢量。
\end{challenge}

\step{模型的边界}

每个模型都有它的管辖范围，越界使用就会闹笑话。本章的边界清单如下。

马吕斯定律假设入射光是完全线偏振光、偏振片是理想的。真实偏振片会吸收掉一部分平行分量，自然光也只是统计意义上的“各方向均等”。拿它算 \(I_0/8\) 时，心里要留着这些零头。

栅栏比喻要明确退场。它\textbf{像的部分}：对绳波这类机械横波，栅栏缝确实只放行平行于缝的振动。它\textbf{不像的部分}：偏振片里没有缝，它的“轴”来自分子层面的有序排列对电场分量的选择性吸收；更致命的是，栅栏只能删不能改，预言“三片栅栏更黑”，而三片偏振片更亮。凡是只用“筛”这个动词思考偏振片的地方，直觉都会出错。

瑞利散射的 \(1/\lambda^4\) 只在颗粒远小于波长时成立。直觉容易顺着“蓝光散射凶”外推：云也是小水滴散射光，那云为什么不是蓝的？因为云滴的直径有十几微米，比波长大几十倍，\(1/\lambda^4\) 的推导前提整个垮掉；大颗粒对各种颜色一视同仁地撒，于是云是白的。同一条“散射”，颗粒尺度换了档，结论就翻了面——这是一个会让直觉出错的反例，值得记住。

色散的“频率越接近固有频率、折射率越大”也有边界：真正撞上固有频率时，推迟不再温和，吸收占据主导，玻璃在那个频率上根本不透明。所谓“反常色散”区域，折射率随频率的行为会把简单的彩虹直觉彻底打乱，那里同时是强吸收带。

双折射只存在于内部有方向性的物质里。普通玻璃各向同性，一个折射率走天下；但把塑料尺子用力掰弯，内部应力会临时造出方向性，把它放在两片偏振镜之间，能看到彩色的应力条纹——工程师正是用这个办法给透明模型“拍受力照片”的。

还有一个购物层面的误用值得单拎出来：偏振镜不等于墨镜。普通深色墨镜只是把各种光一律调暗，对偏振毫无选择；只有标着偏光的镜片才装了那台“投影机器”。区分方法你已经会了：拿镜片去看液晶屏，旋转时有明显明暗变化的是偏振镜，纹丝不动的只是深色玻璃。两者都能挡太阳，但只有偏振镜能消灭水面和路面的眩光——因为那部分光不是“太亮”，而是“太整齐”。

最后，“偏振方向”不要想象成光子飞行途中插着的一根小箭头。它是电磁场振荡方向的整体属性；至于“单个光子的偏振”是什么意思，那要动用概率，是量子篇的问题。

\step{回头解释开场现象}

现在把开场的四件事逐一销账。

三片偏振片复活光线：第一片把自然光投影成竖直线偏振（\(I_0/2\)），中间片把它投影成斜向（\(I_0/4\)），第三片再投影成水平（\(I_0/8\)）。每一片都在“改写”偏振方向而不只是“删”，所以多插一片反而有光。你当初猜的三种说法里，第三种对了：偏振是光身上独立的属性，偏振片是读写它的机器。

手机屏幕旋转镜片变黑：液晶屏的每个像素靠背光照明，背光先过一片偏振片，再经过液晶层——液晶分子的排列受电压控制，能像三片实验里的“中间片”一样旋转偏振方向，最后过第二片偏振片决定这个像素亮还是暗。你的手机屏幕本质上就是几十万个微型“三片偏振片装置”，只不过中间那片换成了电控液晶。这就是为什么屏幕光天生带偏振，也是整个液晶显示工业的原理。

水面反光消失：阳光在水面以约 \(53^\circ\) 的布儒斯特角入射时，反射光几乎是纯的水平线偏振。偏振太阳镜的透光轴做成竖直方向，正好把这份反光投影到零，同时放行一多半来自水下、偏振杂乱的光。钓鱼的人能看清水里的鱼，摄影师用它压掉橱窗反光，都是同一道题。

蓝天与红日落：阳光穿过大气，蓝紫光被分子以四倍于红光的效率撒向全天，你的眼睛接收到的是这份“撒出来的光”，所以天是蓝的；傍晚光路变长几十倍，蓝光在到达你之前几乎撒光，直射而来的太阳只剩橙红。同一个机制，一个方向给你蓝天，另一个方向给你红日——方向反转检查通过。

顺便把两件没排进开场的事也销账。3D 电影：两台放映机各发一种旋转方向的圆偏振，眼镜两片各放行一种，歪头不串画面，原理就是“旋转方向互不混淆”。方解石的双影：晶体给两个偏振方向报出两个折射率，两个方向的光走两条不同的折射路径，一个字就成了两个。至于天空中那片旋转镜片会变暗的天区——那是被撒向你的蓝光恰好以九十度角拐过来，按横波的规矩，这个方向的散射光几乎只剩一个偏振方向，所以能被镜片压暗。蜜蜂每天都在读这张地图，你现在也会读了。

\step{脑洞实验与挑战题}

\begin{thought}[一次只放一个光子]
把光源调暗到极致：平均每一小段时间里只发出一个光子（现代实验室用单光子源早已能做到）。让这样的光子一个接一个地射向两片夹角 \(\theta\) 的偏振片。你会发现：每个光子要么整个穿过、要么整个消失，从来没有“穿过八分之一个光子”。马吕斯定律此时变成了一条概率定律：单个光子穿过的概率是 \(\cos^2\theta\)，亿万光子累计出的强度才是 \(I_0\cos^2\theta\)。

再想想三片装置：斜 \(45^\circ\) 插入中间片后，单个光子的命运被一次次“重新抽签”。偏振方向在这里不再像一根经典的小箭头，而更像一份“遇到下一片时按什么概率抽签”的档案。这个实验是通往第 43 章和第 45 章最平坦的桥：量子态、概率幅与测量，全都能在三片几块钱的塑料片上预演。届时请记住本章的教训——仪器（偏振片）并没有“粗略干扰”光子，概率不是仪器粗糙造成的，它是自然本身的记账方式。
\end{thought}

\begin{puzzles}
\item 两片理想偏振片夹角从 \(0^\circ\) 慢慢转到 \(90^\circ\)，透射光强从 \(I_0\) 一路降到 \(0\)。现在改成三片，第一片与第三片固定垂直，中间片从 \(0^\circ\) 转到 \(90^\circ\)。透射光强是单调变化，还是先升后降？最亮时中间片在什么角度，光强是多少？（提示：把两段投影的 \(\cos^2\) 都写出来，注意两个夹角之和始终是 \(90^\circ\)；试着在 \(0^\circ\)、\(45^\circ\)、\(90^\circ\) 三个点先算再猜。）
\item 雨后夜晚开车，湿漉漉的柏油路面反射车灯特别刺眼。同车的乘客说“戴上偏振墨镜就好了”，司机说“不行，会看不清路面”。谁更有道理？（提示：路面反射的眩光偏振方向如何？从路面缓慢反射回来的、真正携带路面信息的散射光偏振又如何？再想一层：如果所有司机都戴，会少掉什么信息？）
\item 宇航员在大气层外看天空，看到的不是蓝色而是黑色；月球上“天空”同样是黑的，地球的天却是蓝的。请用本章模型解释这中间的差别，并预言：火星大气稀薄且富含尘埃颗粒，它的天空颜色会与地球有何不同？（提示：蓝天需要两个条件——有东西可撒，且撒的家伙比波长小。逐条检查三个天体各自缺了哪条。）
\end{puzzles}
